% Pista ana-23s-q4 · Anàlisi
% Procedència: PAU Catalunya 2023 · Convocatòria de setembre · Sèrie 2 · Qüestió 4
\documentclass[11pt,a4paper]{article}
\usepackage{pista}
\pistainfo{Catalunya 2023}{Convocatòria de setembre}{Sèrie 2}

\begin{document}

\section*{\textcolor{blauPista}{Qüestió 4}}

\noindent Sigui la funció $f(x)$ definida per $f(x) = -3x + e^{2x^{3} - 1}$.

\subsection*{\textit{a)} \normalfont Justifiqueu que $f(x) = 2$ té una solució en l'interval $(-1, 0)$. \hfill\textnormal{[1,25 punts]}}

\pista{Aquest és un problema d'aplicació del Teorema de Bolzano. Però aquest teorema el fem servir per estudiar equacions del tipus $g(x)=0$ i el nostre enunciat diu $f(x)=2$. Una estratègia usual és definir una funció auxiliar $g(x) = f(x) - 2$, perquè llavors ``$f(x) = 2$'' i ``$g(x) = 0$'' són exactament la mateixa condició. Després, comprova que $g$ és contínua a $[-1, 0]$ (mira els seus components: una funció lineal i una altra, exponencial, totes elles contínues) i avalua $g(-1)$ i $g(0)$. Si tenen signes oposats, el teorema et garanteix que existeix un punt dins l'interval on $g(x)$ s'anul·la, és a dir, on $f(x)$ val 2.}

\subsection*{\textit{b)} \normalfont Sigui la funció $h(x) = -3x^{2} + e^{2x^{3} - 1}$. Calculeu l'àrea de la regió compresa entre les gràfiques de les funcions $f(x)$ i $h(x)$. \hfill\textnormal{[1,25 punts]}}

\pista{L'àrea entre dues corbes és la integral definida de la diferència de les dues funcions (la de ``dalt'' menys la de ``baix''). Però abans, calcula la diferència $f(x) - h(x)$ i observa-la amb atenció: hi descobriràs una simplificació molt agradable que farà que la integral sigui senzilla. Per als límits d'integració, has de trobar els punts de tall, és a dir, els valors de $x$ on $f(x) = h(x)$ (que és el mateix que $f(x) - h(x) = 0$, així que ja els tindràs gairebé regalats de la simplificació anterior). Per decidir quina funció va per ``dalt'' a l'interval, avalua la diferència en un punt interior. Ara bé, recorda que també pots calcular el valor absolut de la integral definida, i llavors ja no cal descobrir quina és la funció que queda per ``dalt''.}

\end{document}
